%% ============================================================
%%  SECCIÓN 11 — Análisis de Gráficas
%% ============================================================

\section{Análisis de Gráficas}
\label{sec:analisis_graficas}

El análisis cuantitativo y físico de las representaciones gráficas de la Sección~\ref{sec:graficas} proporciona conclusiones fundamentales sobre el comportamiento del sistema de medición:

\subsection{Análisis de la Curva de Calibración DC (Figura~\ref{fig:calibracion_dc})}
La regresión lineal entre la tensión medida por el osciloscopio en escala óptima ($V_{\text{osc}}$) y la tensión de referencia del multímetro patrón ($V_{\text{mult}}$) demuestra una correlación lineal extraordinaria:
\begin{enumerate}
    \item \textbf{Pendiente del Ajuste ($m = 1{,}0036 \pm 0{,}0039$, con error estándar $s_m = 0{,}0039$):}
    La pendiente experimental concuerda con el valor teórico ideal unitario ($m_{\text{teo}} = 1{,}0000$) dentro de su intervalo expandido a un nivel de confianza del \qty{95}{\percent} ($m \pm 2 s_m = 1{,}0036 \pm 0{,}0078$). El error relativo en la ganancia es de apenas un \qty{0,36}{\percent}, ubicándose holgadamente por debajo de la tolerancia máxima de calibración vertical del fabricante ($\pm\qty{3}{\percent}$ de fondo de escala especificado en~\cite{rs_rtc1000_specs}).
    
    \item \textbf{Ordenada en el Origen ($b = -0{,}021 \pm 0{,}026\,\unit{\volt}$, con error estándar $s_b = \qty{0,026}{\volt}$):}
    El intercepto con el eje vertical es estadísticamente compatible con cero puesto que el valor nulo se ubica dentro de su intervalo de incertidumbre estándar ($|b|/s_b = 0{,}81 < 1$), confirmando que el procedimiento de ajuste del nivel de tierra ($\text{GND}$) eliminó de manera efectiva cualquier tensión residual de \textit{offset} o desbalance en los amplificadores de entrada.

    \item \textbf{Coeficiente de Determinación ($R^2 = 0{,}99995$):}
    El valor de $R^2$ prácticamente unitario ratifica que más del \qty{99,99}{\percent} de la varianza observada es atribuible a la relación lineal estricta entre ambos instrumentos.

    \item \textbf{Distribución de Residuos:}
    En el panel inferior de la Figura~\ref{fig:calibracion_dc}, los residuos se distribuyen aleatoriamente alrededor del valor cero dentro de una banda estrecha inferior a $\pm\qty{0,05}{\volt}$, sin evidenciar patrones no lineales sistemáticos. No obstante, el número de observaciones ($n = 5$, registradas en una sola secuencia ascendente) no permite evaluar de forma concluyente la presencia de histéresis o deriva térmica temporal, para lo cual se requerirían ciclos completos de carga/descarga y mediciones repetidas en el tiempo.
\end{enumerate}

\subsection{Análisis de las Formas de Onda AC (Figura~\ref{fig:ondas_ac})}
\begin{enumerate}
    \item En la onda senoidal del transformador (Figura~\ref{fig:ondas_ac}a), la curva graficada representa la señal senoidal pura de referencia teórica ($V_{p,\text{ref}} = \sqrt{2} V_{\text{rms}} = \qty{9,66}{\volt}$) deducida de la medición patrón True-RMS ($V_{\text{rms}} = \qty{6,830 \pm 0,010}{\volt}$). Esta referencia se contrasta directamente con la amplitud medida en la retícula del osciloscopio ($V_p = \qty{9,60 \pm 0,20}{\volt}$), evidenciando una concordancia extraordinaria con un error relativo de apenas \qty{0,61}{\percent}, contenido holgadamente dentro de la banda de incertidumbre instrumental ($\pm\qty{0,20}{\volt}$). Asimismo, la simetría entre semiciclos arroja un factor de forma $k_f = V_{\text{rms}} / V_{\text{prom,rect}} = \pi / (2\sqrt{2}) \approx 1{,}11$, característico de una oscilación armónica pura de la red comercial de \qty{60}{\hertz}.
    \item En las señales generadas por síntesis digital (Figuras~\ref{fig:ondas_ac}b y~\ref{fig:ondas_ac}c), los valores eficaces calculados ($V_{\text{rms}} = \qty{3,39}{\volt}$ para seno, \qty{2,77}{\volt} para triángulo y \qty{4,80}{\volt} para cuadrada) reproducen con exactitud matemática las relaciones teóricas $1/\sqrt{2}$, $1/\sqrt{3}$ y $1$, con discrepancias inferiores al \qty{0,8}{\percent} respecto a la medición de True-RMS.
\end{enumerate}

\subsection{Análisis del Modo Graficador X-Y y Figuras de Lissajous (Figura~\ref{fig:lissajous_xy})}
\begin{enumerate}
    \item La correspondencia entre el número de lóbulos/tangencias y la relación de frecuencias satisface rigurosamente la ecuación~\eqref{eq:lissajous_ratio}:
    \begin{align}
        \frac{f_y}{f_x} &= \frac{\qty{60}{\hertz}}{\qty{60}{\hertz}} = 1 \implies 1\text{ tangencia horizontal},\; 1\text{ vertical (elipse estacionaria)} \\
        \frac{f_y}{f_x} &= \frac{\qty{120}{\hertz}}{\qty{60}{\hertz}} = 2 \implies 2\text{ tangencias horizontales},\; 1\text{ vertical (figura en ocho)} \\
        \frac{f_y}{f_x} &= \frac{\qty{180}{\hertz}}{\qty{60}{\hertz}} = 3 \implies 3\text{ tangencias horizontales},\; 1\text{ vertical (tres lóbulos)} \\
        \frac{f_y}{f_x} &= \frac{\qty{240}{\hertz}}{\qty{60}{\hertz}} = 4 \implies 4\text{ tangencias horizontales},\; 1\text{ vertical (cuatro lóbulos)}
    \end{align}
    \item La inversión del canal vertical ($\text{CH2 INVERT}$) produce una reflexión especular respecto al eje horizontal $X$ tal que $(x, y) \to (x, -y)$, invirtiendo la inclinación de la elipse sin alterar su excentricidad ni el área encerrada.
    \item En señales DC constantes, la deflexión temporal nula colapsa el haz en un punto estático coordenado en $(+1{,}5\,\unit{\volt}, +3{,}0\,\unit{\volt})$, el cual se refleja simétricamente a $(+1{,}5\,\unit{\volt}, -3{,}0\,\unit{\volt})$ al activar $\text{CH2 INVERT}$, demostrando que el osciloscopio en modo $X$-$Y$ funciona como un trazador de coordenadas cartesianas analógicas directas.
\end{enumerate}
